\section{Conclusion}
\label{sec:conclusion}

We implemented BPR matrix factorization end to end: implicit-feedback
construction, iterative core filtering, point-in-time feature engineering,
vectorized triplet SGD, full-catalogue ranking evaluation, and a
parity-tested serving path deployed by Terraform against a local AWS
emulator.

The measured outcome is that a tuned item-$k$NN baseline edges out tuned BPR
on this data (NDCG@10 $0.1305$ versus $0.1248$). We report it because the
alternative --- tuning until the flagship model wins --- is the specific
failure mode that Dacrema et al.~\cite{dacrema2019} documented across a
large body of published recommender work, and avoiding it was an explicit
goal of this project.

\subsection{What transfers}

The reusable artifacts of this project are not the model weights:

\begin{itemize}
  \item \textbf{A leakage harness.} Future-event injection, fit-partition
        isolation, and label destruction, as executable tests rather than
        prose. The injection test surfaced a genuine dtype-promotion defect
        in the feature layer.
  \item \textbf{A parity harness.} Asserting that the served ranking equals
        the evaluated ranking closes the gap where good offline numbers
        quietly fail to reach production.
  \item \textbf{An honest selection protocol.} Validation-only tuning, with
        the same grid budget extended to the baselines.
\end{itemize}

\subsection{Future work}

\paragraph{Popularity-weighted sampling.} Drawing negatives
$j$ with probability proportional to $c_j^{\alpha}$ rather than uniformly
produces harder triplets as training progresses, and is the most likely
single change to close the observed gap.

\paragraph{Confidence intervals.} Repeating the comparison across seeds and
reporting intervals would establish whether the $0.006$ NDCG gap is
significant.

\paragraph{Hybrid scoring.} The content features in
Table~\ref{tab:features} are currently consumed only for cold-start fallback.
Folding them into the scorer, $\hat{x}_{ui} = \pvec_u^{\top}\qvec_i + b_i +
\mathbf{w}^{\top}\boldsymbol{\phi}_{ui}$, would let the model exploit genre
affinity directly.

\paragraph{Real timestamps.} Evaluating on a dataset with genuine event times
would replace the chronology proxy and make the temporal split defensible as
a simulation of production.
